\editionPage[\pdfbookmark[0]{Indications des défis}{bm-indices-debut-parent}]{indices-debut}{Indications progressives}{Défis 1 à 3}{Ne lire que l’aide nécessaire.}
\begin{blocEdition}{indices-debut}
\editionEtape{D01 · Compter jusqu’au milliard}\textbf{Indice 1.} Combien de secondes dans une journée ? Commencer par $24\times3\,600$.\par\medskip\textbf{Indice 2.} Un milliard, c’est mille millions : le temps est mille fois plus long.\par
\editionEtape{D02 · Entre 3,7 et 3,8}\textbf{Indice 1.} Écrire les bornes $3{,}70$ et $3{,}80$, puis $3{,}700$ et $3{,}800$.\par\medskip\textbf{Indice 2.} Entre deux décimaux, on peut toujours intercaler un autre décimal.\par
\editionEtape{D03 · Somme 5, produit 6,09}\textbf{Indice 1.} Tester $2{,}5$ et $2{,}5$, puis 2 et 3 : que deviennent les produits ?\par\medskip\textbf{Indice 2.} Le produit diminue quand les deux nombres s’éloignent l’un de l’autre.\par

\end{blocEdition}
\editionPage{indices-4}{Indications progressives}{Défis 4 à 6}{Ne lire que l’aide nécessaire.}
\begin{blocEdition}{indices-4}
\editionEtape{D04 · Les restes en escalier}\textbf{Indice 1.} Que se passe-t-il si l’on ajoute 1 au nombre cherché ?\par\medskip\textbf{Indice 2.} Avec 1 de plus, le nombre est un multiple commun de 2, 3, 4, 5 et 6.\par
\editionEtape{D05 · Les pains égyptiens}\textbf{Indice 1.} Commencer par donner à chacun la plus grosse part possible : un demi-pain.\par\medskip\textbf{Indice 2.} Écrire $\frac58$ comme une somme de deux fractions de numérateur 1.\par
\editionEtape{D06 · Une somme qui se simplifie}\textbf{Indice 1.} Les trois différences ont toutes pour numérateur 1.\par\medskip\textbf{Indice 2.} Écrire $\frac12$ comme $1-\frac12$ : beaucoup de termes s’annulent.\par

\end{blocEdition}
\editionPage{indices-7}{Indications progressives}{Défis 7 à 9}{Ne lire que l’aide nécessaire.}
\begin{blocEdition}{indices-7}
\editionEtape{D07 · Hausse puis baisse}\textbf{Indice 1.} Calculer 10~\% du nouveau prix, pas de l’ancien.\par\medskip\textbf{Indice 2.} Les deux pourcentages ne s’appliquent pas au même prix.\par
\editionEtape{D08 · Les tables du banquet}\textbf{Indice 1.} Dessiner les trois premières étapes et repérer les places en bout de rangée.\par\medskip\textbf{Indice 2.} Deux personnes en bout, puis 2 places par table.\par
\editionEtape{D09 · Les engrenages}\textbf{Indice 1.} À chaque tour, la petite roue fait passer 12 dents.\par\medskip\textbf{Indice 2.} Le nombre de dents qui passent est le même pour deux roues en contact.\par

\end{blocEdition}
\editionPage{indices-10}{Indications progressives}{Défis 10 à 12}{Ne lire que l’aide nécessaire.}
\begin{blocEdition}{indices-10}
\editionEtape{D10 · Une ficelle autour de la Terre}\textbf{Indice 1.} Le périmètre vaut $\pi\times d$. De combien augmente le diamètre si le périmètre augmente de 1 m ?\par\medskip\textbf{Indice 2.} Le jeu entre la ficelle et le sol vaut la moitié de l’augmentation du diamètre.\par
\editionEtape{D11 · Le carré penché}\textbf{Indice 1.} Le grand carré se décompose en le carré penché et quatre triangles rectangles.\par\medskip\textbf{Indice 2.} Deux de ces triangles forment un rectangle de 1 cm sur 3 cm.\par
\editionEtape{D12 · Les deux feux}\textbf{Indice 1.} Convertir les deux durées en secondes.\par\medskip\textbf{Indice 2.} Chercher le premier multiple commun de 45 et de 80.\par

\end{blocEdition}
\editionPage{indices-13}{Indications progressives}{Défis 13 à 15}{Ne lire que l’aide nécessaire.}
\begin{blocEdition}{indices-13}
\editionEtape{D13 · Combien de points ?}\textbf{Indice 1.} Tracer le cercle de centre A de rayon 4 cm et faire varier le cercle de centre B.\par\medskip\textbf{Indice 2.} Les deux cercles se touchent quand $r=6-4$ ou $r=6+4$.\par
\editionEtape{D14 · Le centre perdu}\textbf{Indice 1.} Tracer deux cordes du cercle, non parallèles.\par\medskip\textbf{Indice 2.} Le centre est à égale distance des deux extrémités d’une corde.\par
\editionEtape{D15 · Les aiguilles en mouvement}\textbf{Indice 1.} En une minute, la petite aiguille tourne de $30:60=0{,}5^\circ$.\par\medskip\textbf{Indice 2.} Repérer chaque aiguille par son angle à partir de midi.\par

\end{blocEdition}
\editionPage{indices-16}{Indications progressives}{Défis 16 à 18}{Ne lire que l’aide nécessaire.}
\begin{blocEdition}{indices-16}
\editionEtape{D16 · Des triangles d’allumettes}\textbf{Indice 1.} Chercher trois nombres entiers dont la somme vaut 12, rangés du plus petit au plus grand.\par\medskip\textbf{Indice 2.} Le plus grand côté doit être strictement plus petit que la somme des deux autres.\par
\editionEtape{D17 · Deux miroirs parallèles}\textbf{Indice 1.} Travailler sur quadrillage, avec des droites verticales.\par\medskip\textbf{Indice 2.} Compter les carreaux de M à $(d_1)$, puis de $M_1$ à $(d_2)$.\par
\editionEtape{D18 · La pyramide de cubes}\textbf{Indice 1.} Chaque étage est un carré : $1\times1$, $2\times2$, $3\times3$…\par\medskip\textbf{Indice 2.} Ajouter les étages un à un jusqu’à dépasser 200.\par

\end{blocEdition}
\editionPage{indices-19}{Indications progressives}{Défis 19 à 21}{Ne lire que l’aide nécessaire.}
\begin{blocEdition}{indices-19}
\editionEtape{D19 · Le tableau à reconstituer}\textbf{Indice 1.} Commencer par le nombre de garçons, puis celui des filles qui déjeunent à la cantine.\par\medskip\textbf{Indice 2.} Contrôler que chaque ligne et chaque colonne donnent le bon total.\par
\editionEtape{D20 · Les chaussettes dans le noir}\textbf{Indice 1.} Imaginer la pire malchance possible.\par\medskip\textbf{Indice 2.} Avec 2 chaussettes, on peut avoir une rouge et une bleue.\par
\editionEtape{D21 · La spirale du robot}\textbf{Indice 1.} Écrire la liste des longueurs : 1, 1, 2, 2, 3, 3…\par\medskip\textbf{Indice 2.} Suivre la position case par case, ou séparer les déplacements horizontaux et verticaux.\par

\end{blocEdition}
